\documentclass[10pt,a4paper]{article}
\usepackage[margin=2.2cm]{geometry}
\usepackage{fontspec}
\setmonofont{DejaVu Sans Mono}[Scale=0.84]
\newfontfamily\unifont{DejaVu Serif}[Scale=0.9]
\usepackage{amsmath,amssymb,mathtools}
\setcounter{MaxMatrixCols}{20}
\usepackage{longtable,booktabs,array}
\usepackage{fancyvrb}
\usepackage[dvipsnames]{xcolor}
\usepackage[colorlinks=true,linkcolor=NavyBlue,urlcolor=NavyBlue]{hyperref}
\usepackage{microtype}
\usepackage{needspace}
\usepackage{tocloft}\setlength{\cftbeforesecskip}{0.45em}
\setlength{\parindent}{0pt}\setlength{\parskip}{0.55em}
\setlength{\emergencystretch}{8em}
\renewcommand{\arraystretch}{1.15}
\allowdisplaybreaks
\begin{document}
\title{The self-coupled dirac16complex and its Kohn–Sham theory in gravitational fields\\[0.4em]\large The nine-channel self-coupling, the 9×9 Fierz matrix, the covariant nonlinear field equation, chirality in a gravitational field, and Kohn–Sham matter in the primordial, modern and far-future epochs with its back-reaction on the expansion — sections 15–16 of the notebook 7+1.nb}
\author{Repository: github.com/once-ere/7plus1 — prepared for Nash\\ Notebook: 7+1.nb (Wolfram Language 15.0.1), evaluated Tue 27 Oct 2026 20:36:01\\ Document generated by tools/make\_docs.py on 2026-10-27}
\date{}
\maketitle
\begin{abstract}
The 16-component field dirac16complex lives on the 7+1 spacetime (one time coordinate $x_0$, seven space coordinates). This document builds the most general contact self-coupling of the field without derivatives — nine couplings, one for each rank of the Clifford algebra — computes the Fierz matrix that turns exchange terms into direct ones, writes the nonlinear field equation in a gravitational field with the canonical spin connection, and solves the Kohn–Sham theory of homogeneous matter of dirac16complex in the three cosmological epochs, including the self-consistent back-reaction on the expansion through the Einstein–Lovelock equations. Sections 15–16 of 7+1.nb contain 24 checks (24 True); the independent Python program tools/ref\_ks\_gravity.py contains 29 (29 True).
\end{abstract}
\tableofcontents
\Needspace{12\baselineskip}\section*{1. Results at a glance}\addcontentsline{toc}{section}{1. Results at a glance}
\begin{itemize}
\item \textbf{Nine channels.} With $\Theta^K=i^{\sigma_r}\Gamma^{A_1}\cdots\Gamma^{A_r}$ for the 256 increasing multi-indices $K$ the most general self-coupling without derivatives is $\mathcal L_{\rm self}=-\tfrac12\sum_{r=0}^{8}g_rI_r$, $I_r=\sum_{|K|=r}(\bar\psi\Theta_K\psi)(\bar\psi\Theta^K\psi)$ (section 3).
\item \textbf{Fierz.} $F_{rs}=\frac1{16}\sum_l(-1)^{rs-l}\binom{s}{l}\binom{8-s}{r-l}$, $F^2=1$; Hartree–Fock is Hartree with the couplings $G=(1-F^{\mathsf T})g$ (section 4).
\item \textbf{Field equation.} $(i\Gamma^\mu(X)D_\mu-m)\psi=\sum_rg_r\sum_{|K|=r}(\bar\psi\Theta_K\psi)\Theta^K\psi$ with $D_\mu=\partial_\mu+\Omega_\mu$; gravity does not mix the chiral halves (section 5).
\item \textbf{Kohn–Sham matter} in the cosmological spacetime: effective mass $M=m+G_S\rho_S$, vector potential $G_V\rho$; the energy density and the two pressures are integrals over the comoving Fermi ball, and $T^\mu{}_\nu$ is conserved for every expansion history (section 6).
\item \textbf{Epochs.} Near $x_0=0$ the effective mass falls toward 0 (chiral symmetry effectively restored); in the modern epoch the matter is pressureless in ordinary space; in a de Sitter future $1-M$ decays like the density (section 7).
\item \textbf{Back-reaction.} Einstein and Einstein–Lovelock gravity driven by the Kohn–Sham matter: the constraint is preserved by the evolution to a relative $2.4\times10^{-14}$ (section 8).
\end{itemize}
\Needspace{12\baselineskip}\section*{2. Conventions and objects used}\addcontentsline{toc}{section}{2. Conventions and objects used}
Units $\hbar=c=1$; coordinates $X=(x_0,\dots,x_7)$ ($x_0$ time, $x_1,x_2,x_3$ ordinary space, $x_4,\dots,x_7$ auxiliary space); frame metric $\eta=\mathrm{diag}(+1,-1,\dots,-1)$, simply the squares of the octonion basis; repeated indices are summed.

\textbf{The signature: $\{7,1,0\}$ and $\mathrm{diag}(+1,-1,\dots,-1)$.} The 7+1 spacetime has one time-like dimension ($x_0$, the direction of the octonion unit $e_0=1$) and seven space-like dimensions ($x_1,x_2,x_3$ ordinary space, $x_4,\dots,x_7$ auxiliary space). Two different objects describe this, and neither may be replaced by the other. (i) The tensor package xAct, used in section 12 of 7+1.nb exactly as in the author’s notebook Einstein-Lovelock-7+1-Nash.nb, declares the metric with \texttt{DefMetric[\{7, 1, 0\}, g71[-a, -b], CD, …]}: the list $\{7,1,0\}$ declares seven space-like dimensions, one time-like dimension and no degenerate dimension — the numbers of entries $+1$ and $-1$ of $\eta_{71}=\mathrm{diag}(-1,+1,\dots,+1)$ (\texttt{η71}), the same spacetime written in the sign convention of 4+4.nb (space-like $=+1$, time-like $=-1$). (ii) The matrix $\eta=-\eta_{71}=\mathrm{diag}(+1,-1,\dots,-1)$ (\texttt{η17}) is simply the squares of the octonion basis: $e_0^2=+1$ for the unit $e_0=1$ (time) and $e_i^2=-1$ for the seven imaginary units $e_1,\dots,e_7$ (space); more generally it is the real part of the product of two basis units, and it gives the quadratic form of the octonions:
\begin{gather*}
\eta_{AB}=\mathrm{Re}(e_Ae_B),\qquad e_0^2=+1,\qquad e_i^2=-1\ \ (i=1,\dots,7),\\ \mathrm{Re}(x\,x)=x_0^2-x_1^2-\cdots-x_7^2=\eta_{AB}\,x^Ax^B\qquad (x=x^Ae_A\in\mathbb O).
\end{gather*}
$\eta$ gives the signs of the Clifford relation $\Gamma_A\Gamma_B+\Gamma_B\Gamma_A=2\,\eta_{AB}\,1_{16}$ of the real 16×16 Dirac matrices: time squares to $+1_{16}$ and each of the seven space directions to $-1_{16}$, exactly like the octonion units (in the notation $Cl(p,q)$, with $p$ generators squaring to $+1$ and $q$ to $-1$, the algebra $Cl(1,7)\cong\mathbb R(16)$). It also gives the signs of the metric components: $g_{00}>0$ along the time-like direction, negative squares along the space-like directions. Every component computed by 7+1.nb follows from the component matrix of the metric; the declared list fixes for xAct the sign of the determinant of the metric, \texttt{SignDetOfMetric[g71]} $=-1$, which is the same in both sign conventions, since $\det \eta=\det \eta_{71}=-1$. 7+1.nb checks \texttt{SignatureOfMetric[g71]} $=\{7,1,0\}$, \texttt{SignDetOfMetric[g71]} $=-1$, $\mathrm{Re}(e_Ae_B)=\eta_{AB}$ and $\mathrm{Re}(x\,x)=\eta_{AB}x^Ax^B$.

The real 16×16 Dirac matrices and their properties:
\begin{equation*}
\Gamma_0=\begin{pmatrix}0&1_8\\1_8&0\end{pmatrix},\quad \Gamma_i=\begin{pmatrix}\tau^O_i&0\\0&-\tau^O_i\end{pmatrix},\quad \Gamma_A\Gamma_B+\Gamma_B\Gamma_A=2\eta_{AB}1_{16},\quad \Gamma^A=\eta^{AA}\Gamma_A,\quad \bar\psi=\psi^\dagger\Gamma_0,
\end{equation*}
with real antisymmetric 8×8 matrices $\tau^O_i$ ($i=1,\dots,7$; $(\tau^O_i)^2=-1_8$, mutually anticommuting). The volume element and the chirality are
\begin{equation*}
\omega=\Gamma_0\Gamma_1\cdots\Gamma_7=\begin{pmatrix}0&-1_8\\1_8&0\end{pmatrix},\qquad \omega^2=-1_{16},\qquad \chi=i\,\omega,\qquad \chi^2=1_{16},\qquad P_{L,R}=\tfrac12(1_{16}\mp\chi);
\end{equation*}
there is no real chirality matrix. In a gravitational field (frame $e^A{}_\mu$, inverse frame $e_A{}^\mu$, metric $g_{\mu\nu}=\eta_{AB}e^A{}_\mu e^B{}_\nu$) the generalized frame matrices, the canonical spin connection and the covariant derivative are
\begin{align*}
\Gamma^\mu(X)&=e_A{}^\mu\Gamma^A,\qquad \omega_\mu{}^A{}_B=e^A{}_\nu\big(\partial_\mu e_B{}^\nu+\{{}^{\nu}_{\mu\lambda}\}e_B{}^\lambda\big),\qquad \omega_{\mu AB}=\eta_{AC}\,\omega_\mu{}^C{}_B, \\
\Omega_\mu&=\tfrac14\omega_{\mu AB}\Gamma^A\Gamma^B,\qquad D_\mu\psi=\partial_\mu\psi+\Omega_\mu\psi,\qquad \{{}^{\lambda}_{\mu\nu}\}=\tfrac12g^{\lambda\rho}(\partial_\mu g_{\rho\nu}+\partial_\nu g_{\rho\mu}-\partial_\rho g_{\mu\nu}).
\end{align*}
The cosmological metric used for the epochs, its Hubble rates (a prime is $d/dx_0$) and its spin connection:
\begin{align*}
ds^2&=dx_0^2-a^2(dx_1^2+dx_2^2+dx_3^2)-b^2(dx_4^2+\cdots+dx_7^2),\qquad H_a=a'/a,\ \ H_b=b'/b, \\
\Omega_0&=0,\quad \Omega_i=\tfrac{a'}{2}\Gamma_i\Gamma_0\ (i=1,2,3),\quad \Omega_k=\tfrac{b'}{2}\Gamma_k\Gamma_0\ (k=4,\dots,7),\quad \Gamma^\mu\Omega_\mu=\big(\tfrac32H_a+2H_b\big)\Gamma_0 .
\end{align*}
\Needspace{12\baselineskip}\section*{3. The nine channels of the self-coupling (section 15.1)}\addcontentsline{toc}{section}{3. The nine channels of the self-coupling (section 15.1)}
For a strictly increasing multi-index $K=(A_1<\dots<A_r)$ of rank $r$ ($\Gamma^{()}=1_{16}$):
\begin{equation*}
\Theta^K=i^{\sigma_r}\,\Gamma^{A_1}\cdots\Gamma^{A_r},\quad \sigma_r=\tfrac{r(r-1)}{2}\bmod 2,\qquad \Theta_K=\eta_K\Theta^K,\quad \eta_K=\eta_{A_1A_1}\cdots\eta_{A_rA_r}.
\end{equation*}
The phase makes every bilinear $\bar\psi\Theta^K\psi$ real ($\Gamma_0\Theta^K$ is Hermitian), and $\Theta_K\Theta^K=1_{16}$, $\mathrm{tr}(\Theta_K\Theta^J)=16\delta_{KJ}$. The invariant of rank $r$ and the general self-coupling are
\begin{equation*}
I_r=\sum_{|K|=r}(\bar\psi\Theta_K\psi)(\bar\psi\Theta^K\psi)=\sum_{|K|=r}\eta_K(\bar\psi\Theta^K\psi)^2,\qquad \mathcal L_{\rm self}=-\frac12\sum_{r=0}^{8}g_r\,I_r .
\end{equation*}
The ranks have $1,8,28,56,70,56,28,8,1$ members (256 in all). The earlier Kohn–Sham theory of the notebook (section 8 of 7+1.nb) has four couplings, scalar $g_S$, vector $g_V$, axial $g_A$ and pseudoscalar $g_P$:
\begin{equation*}
\mathcal L_{\rm self}=-\frac{g_S}2(\bar\psi\psi)^2-\frac{g_P}2(\bar\psi i\chi\psi)^2-\frac{g_V}2(\bar\psi\Gamma^A\psi)(\bar\psi\Gamma_A\psi)-\frac{g_A}2(\bar\psi\Gamma^A\chi\psi)(\bar\psi\Gamma_A\chi\psi).
\end{equation*}
They are four of the nine:
\begin{gather*}
g_0=g_S,\quad g_1=g_V,\quad g_7=-g_A,\quad g_8=-g_P;\\ I_0=(\bar\psi\psi)^2,\ \ I_1=(\bar\psi\Gamma^A\psi)(\bar\psi\Gamma_A\psi),\ \ I_8=-(\bar\psi\,i\chi\,\psi)^2,\ \ I_7=-(\bar\psi\Gamma^A\chi\psi)(\bar\psi\Gamma_A\chi\psi).
\end{gather*}
$g_2,\dots,g_6$ are new (tensor channels of ranks 2 to 6). Each $I_r$ is invariant under every local Lorentz transformation $\psi\to S\psi$, $\bar\psi\to\bar\psi S^{-1}$, $S=\exp\big(\tfrac18w_{AB}[\Gamma^A,\Gamma^B]\big)$ with real antisymmetric $w_{AB}$ ($S^{-1}=\Gamma_0S^\dagger\Gamma_0$), so in a gravitational field it is written with the frame matrices exactly as in flat spacetime. $\Theta^K\omega$ is, up to a phase of modulus 1, $\Theta$ of the complementary multi-index (Hodge duality).

\Needspace{12\baselineskip}\section*{4. The Fierz matrix and Hartree–Fock with nine channels (section 15.2)}\addcontentsline{toc}{section}{4. The Fierz matrix and Hartree–Fock with nine channels (section 15.2)}
The completeness of the 256 matrices, $\sum_K(\Theta_K)_{cd}(\Theta^K)_{ab}=16\,\delta_{ad}\delta_{bc}$, rearranges every product of two bilinears. For commuting spinors $a,b,c,d$ (anticommuting fields add a minus sign)
\begin{equation*}
\sum_{K\in r}(a\Theta_Kb)(c\Theta^Kd)=\sum_{s=0}^{8}F_{rs}\sum_{J\in s}(a\Theta_Jd)(c\Theta^Jb),\qquad F_{rs}=\frac1{16}\sum_{l}(-1)^{rs-l}\binom{s}{l}\binom{8-s}{r-l},
\end{equation*}
and $F_{rs}=\frac1{256}\sum_{K\in r}\mathrm{tr}(\Theta_K\Theta_J\Theta^K\Theta^J)$ for any $J$ of rank $s$. The matrix $16F$ (rows $r=0,\dots,8$, columns $s=0,\dots,8$):

{\footnotesize
\begin{longtable}{>{\raggedright\arraybackslash}p{\dimexpr 0.1000\linewidth-2\tabcolsep\relax}>{\raggedright\arraybackslash}p{\dimexpr 0.1000\linewidth-2\tabcolsep\relax}>{\raggedright\arraybackslash}p{\dimexpr 0.1000\linewidth-2\tabcolsep\relax}>{\raggedright\arraybackslash}p{\dimexpr 0.1000\linewidth-2\tabcolsep\relax}>{\raggedright\arraybackslash}p{\dimexpr 0.1000\linewidth-2\tabcolsep\relax}>{\raggedright\arraybackslash}p{\dimexpr 0.1000\linewidth-2\tabcolsep\relax}>{\raggedright\arraybackslash}p{\dimexpr 0.1000\linewidth-2\tabcolsep\relax}>{\raggedright\arraybackslash}p{\dimexpr 0.1000\linewidth-2\tabcolsep\relax}>{\raggedright\arraybackslash}p{\dimexpr 0.1000\linewidth-2\tabcolsep\relax}>{\raggedright\arraybackslash}p{\dimexpr 0.1000\linewidth-2\tabcolsep\relax}}
\toprule
\textbf{$r\backslash s$} & \textbf{0} & \textbf{1} & \textbf{2} & \textbf{3} & \textbf{4} & \textbf{5} & \textbf{6} & \textbf{7} & \textbf{8} \\
\midrule\endfirsthead
\toprule
\textbf{$r\backslash s$} & \textbf{0} & \textbf{1} & \textbf{2} & \textbf{3} & \textbf{4} & \textbf{5} & \textbf{6} & \textbf{7} & \textbf{8} \\
\midrule\endhead
\endfoot
\bottomrule\endlastfoot
0 & 1 & 1 & 1 & 1 & 1 & 1 & 1 & 1 & 1 \\
1 & 8 & -6 & 4 & -2 & 0 & 2 & -4 & 6 & -8 \\
2 & 28 & 14 & 4 & -2 & -4 & -2 & 4 & 14 & 28 \\
3 & 56 & -14 & -4 & 6 & 0 & -6 & 4 & 14 & -56 \\
4 & 70 & 0 & -10 & 0 & 6 & 0 & -10 & 0 & 70 \\
5 & 56 & 14 & -4 & -6 & 0 & 6 & 4 & -14 & -56 \\
6 & 28 & -14 & 4 & 2 & -4 & 2 & 4 & -14 & 28 \\
7 & 8 & 6 & 4 & 2 & 0 & -2 & -4 & -6 & -8 \\
8 & 1 & -1 & 1 & -1 & 1 & -1 & 1 & -1 & 1 \\*
\end{longtable}}
(computed here from the closed form; the matrix printed by the notebook is identical). $F^2=1$, with five eigenvalues $+1$ and four $-1$. With a density matrix $R$ the direct (Hartree) energy is $\sum_s\frac{g_s}{2}D_s$, $D_s=\sum_{J\in s}\mathrm{tr}(\Theta_JR)\,\mathrm{tr}(\Theta^JR)$, and the exchange (Fock) energy is $-\sum_r\frac{g_r}{2}\sum_{K\in r}\mathrm{tr}(\Theta_KR\Theta^KR)=-\sum_{r,s}\frac{g_r}{2}F_{rs}D_s$. Hartree–Fock is therefore Hartree with the Fierz-transformed couplings
\begin{equation*}
G=(1-F^{\mathsf T})\,g,\qquad G_s=g_s-\sum_rF_{rs}\,g_r .
\end{equation*}
For uniform spin-saturated matter $R=(\rho_S1_{16}+\rho\,\Gamma_0)/16$ only $D_0=\rho_S^2$ and $D_1=\rho^2$ survive, so only $G_0$ and $G_1$ matter:
\begin{align*}
G_0&=\tfrac{15}{16}g_0-\tfrac{1}{2}g_1-\tfrac{7}{4}g_2-\tfrac{7}{2}g_3-\tfrac{35}{8}g_4-\tfrac{7}{2}g_5-\tfrac{7}{4}g_6-\tfrac{1}{2}g_7-\tfrac{1}{16}g_8, \\
G_1&=-\tfrac{1}{16}g_0+\tfrac{11}{8}g_1-\tfrac{7}{8}g_2+\tfrac{7}{8}g_3-\tfrac{7}{8}g_5+\tfrac{7}{8}g_6-\tfrac{3}{8}g_7+\tfrac{1}{16}g_8 .
\end{align*}
With the couplings of the earlier sections ($g_S=-100000$, $g_V=10000$, all others 0): $G_0=-98750$, $G_1=20000$; an additional tensor coupling $g_2=3000$ gives $(G_0,G_1)=(-104000,17375)$. With $g=(g_S,g_V,0,0,0,0,0,-g_A,-g_P)$ the exchange energy of uniform matter is that of the four-coupling theory,
\begin{equation*}
E_x=-\frac{g_S}{32}(\rho_S^2+\rho^2)+\frac{g_P}{32}(\rho_S^2-\rho^2)-\frac{g_V}{32}(8\rho_S^2-6\rho^2)+\frac{g_A}{32}(8\rho_S^2+6\rho^2).
\end{equation*}
\Needspace{12\baselineskip}\section*{5. The covariant nonlinear field equation (section 15.3) and chirality (section 15.4)}\addcontentsline{toc}{section}{5. The covariant nonlinear field equation (section 15.3) and chirality (section 15.4)}
With $\mathcal L=\sqrt{|g|}\,[\bar\psi(i\Gamma^\mu D_\mu-m)\psi+\mathcal L_{\rm self}]$ the field equation is
\begin{equation*}
\big(i\,\Gamma^\mu(X)\,D_\mu-m\big)\psi=\sum_{r=0}^{8}g_r\sum_{|K|=r}(\bar\psi\Theta_K\psi)\,\Theta^K\psi .
\end{equation*}
The bilinears are Grassmann-even, so the left derivative with respect to $\bar\psi_a$ acts on them as on commuting variables; the check differentiates $\mathcal L_{\rm self}$ with symbolic $\bar\psi$. The chirality $\chi$ commutes with every generator $\frac14[\Gamma^A,\Gamma^B]$, hence with every $\Omega_\mu$: gravity does not mix the chiral halves. With the chiral basis
\begin{gather*}
E_L=\frac1{\sqrt2}\begin{pmatrix}1_8\\-i\,1_8\end{pmatrix},\quad E_R=\frac1{\sqrt2}\begin{pmatrix}-i\,1_8\\1_8\end{pmatrix},\qquad \chi E_L=-E_L,\ \ \chi E_R=E_R,\\ \psi=E_L\phi_L+E_R\phi_R,\ \ \phi_{L,R}=E_{L,R}^\dagger\psi
\end{gather*}
(8-component $\phi_{L,R}$; $E_L^\dagger E_L=E_R^\dagger E_R=1_8$, $E_L^\dagger E_R=0$) and
\begin{equation*}
(\tau^O\!\cdot\partial)_{(a,b)}=\frac1a\sum_{i=1}^{3}\tau^O_i\partial_i+\frac1b\sum_{k=4}^{7}\tau^O_k\partial_k,
\end{equation*}
the free Dirac equation in the cosmological spacetime becomes the two coupled 8-component equations
\begin{equation*}
\big(i\partial_0+i(\tfrac32H_a+2H_b)+(\tau^O\!\cdot\partial)_{(a,b)}\big)\phi_L=m\,\phi_R,\qquad \big(i\partial_0+i(\tfrac32H_a+2H_b)-(\tau^O\!\cdot\partial)_{(a,b)}\big)\phi_R=m\,\phi_L,
\end{equation*}
obtained as $E_{L}^\dagger\Gamma_0(i\Gamma^\mu D_\mu-m)\psi$ and $E_{R}^\dagger\Gamma_0(i\Gamma^\mu D_\mu-m)\psi$.

\Needspace{12\baselineskip}\section*{6. Kohn–Sham theory in gravitational fields (section 16)}\addcontentsline{toc}{section}{6. Kohn–Sham theory in gravitational fields (section 16)}
For a STATIC metric the energy of the self-coupled field is a functional of its densities (defined with the frame matrices and integrated with $\sqrt{|g|}\,d^7x$); the ground state minimizes it (Hohenberg–Kohn), and the Kohn–Sham orbitals $\varphi_n$ solve the covariant Dirac equation with the mean fields,
\begin{equation*}
\big(i\Gamma^\mu(X)D_\mu-m-\Sigma(X)\big)\varphi_n=0,\qquad \Sigma=\sum_rg_r\sum_{|K|=r}\big[\mathrm{tr}(\Theta_K\mathcal R)\,\Theta^K-\Theta^K\mathcal R\,\Theta_K\big],\qquad \mathcal R=\sum_{n\ \rm occ}\varphi_n\bar\varphi_n
\end{equation*}
(direct minus exchange; with the Fierz matrix the exchange part is again of the direct form, with $G$). For a time-dependent metric the same equations with the time-dependent $\mathcal R(X)$ are the time-dependent Kohn–Sham equations; their adiabatic limit holds when the expansion rates are small compared with the orbital energies. For homogeneous, spin-saturated matter in the cosmological spacetime the orbitals are the modes $\chi_k=e^{i(\vec p\cdot\vec x+\vec q\cdot\vec y)}c(x_0)$, $\psi=(a^3b^4)^{-1/2}\chi_k$, with comoving momenta $k=(\vec p,\vec q)$ ($\vec p$ along ordinary space $\vec x=(x_1,x_2,x_3)$, $\vec q$ along the auxiliary space $\vec y=(x_4,\dots,x_7)$) inside a comoving Fermi ball $|k|<k_F$, and the mean fields reduce to an effective mass and a vector potential:
\begin{equation*}
i\frac{dc}{dx_0}=\Big(\sum_{i=1}^3\alpha_i\frac{p_i}{a}+\sum_{k=4}^7\alpha_k\frac{q_k}{b}+\beta M(x_0)+V_0(x_0)\Big)c,\qquad M=m+G_S\rho_S,\quad V_0=G_V\rho,
\end{equation*}
with $\alpha_j=\Gamma^0\Gamma^j$, $\beta=\Gamma_0$, and $G_S\equiv G_0$, $G_V\equiv G_1$ the Hartree–Fock couplings of section 4. The physical densities are integrals over the comoving ball with $|k_a|=k\cos\theta$ (ordinary space), $|k_b|=k\sin\theta$ (auxiliary space), the solid-angle weight $\cos^2\theta\sin^3\theta$ and $\omega=\sqrt{k^2\cos^2\theta/a^2+k^2\sin^2\theta/b^2+M^2}$:
\begin{align*}
\langle f\rangle&=\frac{1}{2\pi^4a^3b^4}\int_0^{k_F}dk\,k^6\int_0^{\pi/2}d\theta\,\cos^2\theta\,\sin^3\theta\;f,\qquad \rho=\langle1\rangle=\frac{k_F^7}{105\pi^4a^3b^4},\qquad \rho_S=\Big\langle\frac{M}{\omega}\Big\rangle, \\
\varepsilon&=\langle\omega\rangle-\frac{(M-m)^2}{2G_S}+\frac{G_V}{2}\rho^2,\qquad p_a=\Big\langle\frac{k^2\cos^2\theta}{3a^2\omega}\Big\rangle+\Pi,\qquad p_b=\Big\langle\frac{k^2\sin^2\theta}{4b^2\omega}\Big\rangle+\Pi, \\
\Pi&=\frac{(M-m)^2}{2G_S}+\frac{G_V}{2}\rho^2,\qquad M=m+G_S\,\rho_S(M)\ \ \text{(gap equation)} .
\end{align*}
$T^\mu{}_\nu=\mathrm{diag}(\varepsilon,-p_a,-p_a,-p_a,-p_b,-p_b,-p_b,-p_b)$ is conserved for every expansion history,
\begin{equation*}
\varepsilon'+3H_a(\varepsilon+p_a)+4H_b(\varepsilon+p_b)=0,
\end{equation*}
because $\varepsilon$ is stationary in $M$ (the gap equation), which removes the term in $M'$ (checked for $a=1+0.3x_0+0.05x_0^2$, $b=1/(1+0.2x_0)$ at $x_0=0.7$: residual$/\varepsilon'$ $=-4.0\times10^{-12}$). The integrals are evaluated with Gauss–Legendre quadrature (96 nodes in $k$, 64 in $\theta$). Units and couplings: $m=1$; $G_S=-100000$, $G_V=10000$. The comoving Fermi momentum is the SATURATION value $k_{F,\rm sat}$ of flat spacetime ($a=b=1$): the Fermi momentum at which the energy per particle $E/A=\varepsilon/\rho$ is minimal, equivalently at which the pressure vanishes, $p=\rho\,\varepsilon_F-\varepsilon=0$ (Gibbs–Duhem), with the Fermi energy $\varepsilon_F=G_V\rho+\sqrt{k_F^2+M^2}$. In flat spacetime the notebook obtains
\begin{gather*}
m^*=M=0.348097340146,\qquad \rho=1.388100213\times10^{-5},\qquad \varepsilon=1.348787233\times10^{-5},\\ \frac{E}{A}-m=\frac{\varepsilon}{\rho}-m=-0.02832142796,\qquad p_a=p_b=0 .
\end{gather*}
\Needspace{12\baselineskip}\section*{7. The three epochs (section 16.2)}\addcontentsline{toc}{section}{7. The three epochs (section 16.2)}
\textbf{Primordial} — the Einstein–Kasner background $a=x_0^{p}$, $b=x_0^{q}$ with $p=(1+2\sqrt2)/7$, $q=(2-3\sqrt2)/14$ (normalized to $a=b=1$ at $x_0=1$; it solves the vacuum equations $G^\mu{}_\nu=0$, whose Kasner conditions are $3p+4q=1$, $3p^2+4q^2=1$). The density grows like $1/x_0$ ($a^3b^4=x_0^{3p+4q}=x_0$), but the ordinary momenta are blue-shifted ($p_i/a\to\infty$), so the scalar density per particle $M/\omega$ falls and the attractive scalar coupling drives the effective mass toward zero — chiral symmetry is effectively restored near $x_0=0$:

{\footnotesize
\begin{longtable}{>{\raggedright\arraybackslash}p{\dimexpr 0.1600\linewidth-2\tabcolsep\relax}>{\raggedright\arraybackslash}p{\dimexpr 0.1400\linewidth-2\tabcolsep\relax}>{\raggedright\arraybackslash}p{\dimexpr 0.1400\linewidth-2\tabcolsep\relax}>{\raggedright\arraybackslash}p{\dimexpr 0.1400\linewidth-2\tabcolsep\relax}>{\raggedright\arraybackslash}p{\dimexpr 0.1400\linewidth-2\tabcolsep\relax}>{\raggedright\arraybackslash}p{\dimexpr 0.1400\linewidth-2\tabcolsep\relax}>{\raggedright\arraybackslash}p{\dimexpr 0.1400\linewidth-2\tabcolsep\relax}}
\toprule
\textbf{$x_0$} & \textbf{$10^{-4}$} & \textbf{$10^{-3}$} & \textbf{$10^{-2}$} & \textbf{$10^{-1}$} & \textbf{$1$} & \textbf{$10$} \\
\midrule\endfirsthead
\toprule
\textbf{$x_0$} & \textbf{$10^{-4}$} & \textbf{$10^{-3}$} & \textbf{$10^{-2}$} & \textbf{$10^{-1}$} & \textbf{$1$} & \textbf{$10$} \\
\midrule\endhead
\endfoot
\bottomrule\endlastfoot
$m^*$ & 0.00382418 & 0.0107840 & 0.0302385 & 0.0884446 & 0.348097 & 0.891764 \\*
\end{longtable}}
The pressure is anisotropic: $p_a/p_b=1.08689$, $1.64617$, $-0.403633$ at $x_0=10^{-3}$, $10^{-1}$, $10$ (at $x_0=1$, flat spacetime at the saturation density, both pressures vanish).

\textbf{Modern} — a static auxiliary space ($b=1$) and growing $a$: the ordinary momenta redshift, $M\to m$, $\varepsilon a^3\to$ const and $p_a/\varepsilon\to0$ — dirac16complex behaves as cold (pressureless) matter in ordinary space, while a stress remains along the auxiliary space ($p_b/\varepsilon=0.0497735$ at $a=100$), which the stabilizing stress of the modern epoch must balance:

{\footnotesize
\begin{longtable}{>{\raggedright\arraybackslash}p{\dimexpr 0.1600\linewidth-2\tabcolsep\relax}>{\raggedright\arraybackslash}p{\dimexpr 0.1680\linewidth-2\tabcolsep\relax}>{\raggedright\arraybackslash}p{\dimexpr 0.1680\linewidth-2\tabcolsep\relax}>{\raggedright\arraybackslash}p{\dimexpr 0.1680\linewidth-2\tabcolsep\relax}>{\raggedright\arraybackslash}p{\dimexpr 0.1680\linewidth-2\tabcolsep\relax}>{\raggedright\arraybackslash}p{\dimexpr 0.1680\linewidth-2\tabcolsep\relax}}
\toprule
\textbf{$a$} & \textbf{$1$} & \textbf{$2$} & \textbf{$4$} & \textbf{$10$} & \textbf{$100$} \\
\midrule\endfirsthead
\toprule
\textbf{$a$} & \textbf{$1$} & \textbf{$2$} & \textbf{$4$} & \textbf{$10$} & \textbf{$100$} \\
\midrule\endhead
\endfoot
\bottomrule\endlastfoot
$m^*$ & 0.34809734 & 0.85354081 & 0.98074107 & 0.99875760 & 0.99999876 \\*
\end{longtable}}
\textbf{Far future} — de Sitter expansion of ordinary space $a=e^{0.1x_0}$ over $b=1$: $\rho=\rho_0e^{-0.3x_0}$ and $1-M$ decays like $\rho$:

{\footnotesize
\begin{longtable}{>{\raggedright\arraybackslash}p{\dimexpr 0.1600\linewidth-2\tabcolsep\relax}>{\raggedright\arraybackslash}p{\dimexpr 0.1680\linewidth-2\tabcolsep\relax}>{\raggedright\arraybackslash}p{\dimexpr 0.1680\linewidth-2\tabcolsep\relax}>{\raggedright\arraybackslash}p{\dimexpr 0.1680\linewidth-2\tabcolsep\relax}>{\raggedright\arraybackslash}p{\dimexpr 0.1680\linewidth-2\tabcolsep\relax}>{\raggedright\arraybackslash}p{\dimexpr 0.1680\linewidth-2\tabcolsep\relax}}
\toprule
\textbf{$x_0$} & \textbf{$0$} & \textbf{$10$} & \textbf{$20$} & \textbf{$40$} & \textbf{$80$} \\
\midrule\endfirsthead
\toprule
\textbf{$x_0$} & \textbf{$0$} & \textbf{$10$} & \textbf{$20$} & \textbf{$40$} & \textbf{$80$} \\
\midrule\endhead
\endfoot
\bottomrule\endlastfoot
$m^*$ & 0.34809734 & 0.93957969 & 0.99692358 & 0.99999236 & 1.00000000 \\*
\end{longtable}}
\Needspace{12\baselineskip}\section*{8. Self-consistent evolution: Einstein–Lovelock gravity driven by Kohn–Sham matter (section 16.3)}\addcontentsline{toc}{section}{8. Self-consistent evolution: Einstein–Lovelock gravity driven by Kohn–Sham matter (section 16.3)}
With $T^\mu{}_\nu$ of the Kohn–Sham matter the Einstein–Lovelock field equations
\begin{equation*}
\sum_{k=1}^{3}\alpha_kE^{(k)\mu}{}_\nu-\bar\Lambda\,\delta^\mu_\nu=\kappa\,T^\mu{}_\nu,\qquad \alpha_1=1,
\end{equation*}
with the Lovelock tensors
\begin{equation*}
E^{(k)\mu}{}_\nu=-\frac{1}{2^{k+1}}\,\delta^{\mu\alpha_1\beta_1\cdots\alpha_k\beta_k}_{\nu\gamma_1\delta_1\cdots\gamma_k\delta_k}\,R^{\gamma_1\delta_1}{}_{\alpha_1\beta_1}\cdots R^{\gamma_k\delta_k}{}_{\alpha_k\beta_k},\qquad \delta^{\mu_1\cdots\mu_n}_{\nu_1\cdots\nu_n}=\det\big[\delta^{\mu_i}_{\nu_j}\big],
\end{equation*}
($E^{(1)}$ is the Einstein tensor; $R^{\gamma\delta}{}_{\alpha\beta}$ the Riemann tensor) are a closed system for $a,b$. For the cosmological metric the needed components are ($H_a'=dH_a/dx_0$, $H_b'=dH_b/dx_0$):
\begin{align*}
E^{(1)0}{}_0&=3 \left(H_{a}^{2} + 4 H_{a} H_{b} + 2 H_{b}^{2}\right), \\
E^{(1)1}{}_1&=2 H'_{a} + 4 H'_{b} + 3 H_{a}^{2} + 8 H_{a} H_{b} + 10 H_{b}^{2}, \\
E^{(1)4}{}_4&=3 H'_{a} + 3 H'_{b} + 6 H_{a}^{2} + 9 H_{a} H_{b} + 6 H_{b}^{2},
\end{align*}
\begin{align*}
E^{(2)0}{}_0&=- 12 H_{b} \left(4 H_{a}^{3} + 18 H_{a}^{2} H_{b} + 12 H_{a} H_{b}^{2} + H_{b}^{3}\right), \\
E^{(2)1}{}_1&=- 32 H'_{a} H_{a} H_{b} - 48 H'_{a} H_{b}^{2} - 16 H'_{b} H_{a}^{2} - 96 H'_{b} H_{a} H_{b} \\
&\quad - 48 H'_{b} H_{b}^{2} - 32 H_{a}^{3} H_{b} - 136 H_{a}^{2} H_{b}^{2} - 192 H_{a} H_{b}^{3} \\
&\quad - 60 H_{b}^{4}, \\
E^{(2)4}{}_4&=- 12 H'_{a} H_{a}^{2} - 72 H'_{a} H_{a} H_{b} - 36 H'_{a} H_{b}^{2} - 36 H'_{b} H_{a}^{2} \\
&\quad - 72 H'_{b} H_{a} H_{b} - 12 H'_{b} H_{b}^{2} - 12 H_{a}^{4} - 108 H_{a}^{3} H_{b} \\
&\quad - 180 H_{a}^{2} H_{b}^{2} - 108 H_{a} H_{b}^{3} - 12 H_{b}^{4},
\end{align*}
\begin{align*}
E^{(3)0}{}_0&=360 H_{a}^{2} H_{b}^{3} \left(4 H_{a} + 3 H_{b}\right), \\
E^{(3)1}{}_1&=72 H_{b}^{2} \left(8 H'_{a} H_{a} H_{b} + 2 H'_{a} H_{b}^{2} + 12 H'_{b} H_{a}^{2} + 8 H'_{b} H_{a} H_{b} + 8 H_{a}^{3} H_{b} + 19 H_{a}^{2} H_{b}^{2} + 8 H_{a} H_{b}^{3}\right), \\
E^{(3)4}{}_4&=72 H_{a} H_{b} \left(9 H'_{a} H_{a} H_{b} + 6 H'_{a} H_{b}^{2} + 6 H'_{b} H_{a}^{2} + 9 H'_{b} H_{a} H_{b} + 9 H_{a}^{3} H_{b} + 17 H_{a}^{2} H_{b}^{2} + 9 H_{a} H_{b}^{3}\right).
\end{align*}
The 11 and 44 equations are LINEAR in $(H_a',H_b')$ (Lovelock equations are quasi-linear) and are solved for them at every step; the 00 equation contains no second derivatives and is a CONSTRAINT, imposed on the initial data ($H_b(0)$ from $H_a(0)$) and preserved by the evolution because both sides are conserved. The runs ($\kappa=1000$, $\bar\Lambda=0$, $H_a(0)=0.05$, $a(0)=b(0)=1$, $x_0$ from 0 to 20; explicit Runge–Kutta of order 8 with InterpolationOrder → All, so that the solution between the steps has the order of the method) give:

{\footnotesize
\begin{longtable}{>{\raggedright\arraybackslash}p{\dimexpr 0.2200\linewidth-2\tabcolsep\relax}>{\raggedright\arraybackslash}p{\dimexpr 0.1300\linewidth-2\tabcolsep\relax}>{\raggedright\arraybackslash}p{\dimexpr 0.1300\linewidth-2\tabcolsep\relax}>{\raggedright\arraybackslash}p{\dimexpr 0.1300\linewidth-2\tabcolsep\relax}>{\raggedright\arraybackslash}p{\dimexpr 0.1300\linewidth-2\tabcolsep\relax}>{\raggedright\arraybackslash}p{\dimexpr 0.1300\linewidth-2\tabcolsep\relax}>{\raggedright\arraybackslash}p{\dimexpr 0.1300\linewidth-2\tabcolsep\relax}}
\toprule
\textbf{gravity} & \textbf{$H_b(0)$} & \textbf{$a(20)$} & \textbf{$b(20)$} & \textbf{$H_a(20)$} & \textbf{$H_b(20)$} & \textbf{max. relative constraint violation} \\
\midrule\endfirsthead
\toprule
\textbf{gravity} & \textbf{$H_b(0)$} & \textbf{$a(20)$} & \textbf{$b(20)$} & \textbf{$H_a(20)$} & \textbf{$H_b(20)$} & \textbf{max. relative constraint violation} \\
\midrule\endhead
\endfoot
\bottomrule\endlastfoot
Einstein ($\alpha_2=\alpha_3=0$) & 0.0091437124 & 1.5817387 & 1.1975586 & 0.011877999 & 0.0074191078 & $2.0\times10^{-14}$ \\
Einstein–Lovelock ($\alpha_2=20$, $\alpha_3=500$) & 0.014980704 & 1.4773515 & 1.2760841 & 0.010175914 & 0.0084887647 & $2.4\times10^{-14}$ \\*
\end{longtable}}
\Needspace{12\baselineskip}\section*{9. Verification}\addcontentsline{toc}{section}{9. Verification}
Every statement above is a check of 7+1.nb (Wolfram Language 15.0.1) and of the independent Python program tools/ref\_ks\_gravity.py (numpy, scipy, sympy), which rebuilds the Dirac matrices and computes everything again with its own quadrature, gap solver and Runge–Kutta integrator (DOP853). The numbers printed by both programs:

{\footnotesize
\begin{longtable}{>{\raggedright\arraybackslash}p{\dimexpr 0.2400\linewidth-2\tabcolsep\relax}>{\raggedright\arraybackslash}p{\dimexpr 0.3100\linewidth-2\tabcolsep\relax}>{\raggedright\arraybackslash}p{\dimexpr 0.3100\linewidth-2\tabcolsep\relax}>{\raggedright\arraybackslash}p{\dimexpr 0.1400\linewidth-2\tabcolsep\relax}}
\toprule
\textbf{quantity} & \textbf{7+1.nb} & \textbf{Python} & \textbf{max. relative difference} \\
\midrule\endfirsthead
\toprule
\textbf{quantity} & \textbf{7+1.nb} & \textbf{Python} & \textbf{max. relative difference} \\
\midrule\endhead
\endfoot
\bottomrule\endlastfoot
flat limit: $m^*$, $\rho$, $\varepsilon$ & 0.34809734, $1.3881002\times10^{-5}$, $1.3487872\times10^{-5}$ & 0.34809734, $1.3881002\times10^{-5}$, $1.3487872\times10^{-5}$ & $7.5\times10^{-15}$ \\
primordial $m^*(x_0)$ & 0.0038241794, 0.010783990, 0.030238520, 0.088444641, 0.34809734, 0.89176375 & 0.0038241794, 0.010783990, 0.030238520, 0.088444641, 0.34809734, 0.89176375 & $5.3\times10^{-15}$ \\
primordial $p_a/p_b$ & 1.0868869, 1.6461738, -0.40363339 & 1.0868869, 1.6461738, -0.40363339 & $2.3\times10^{-15}$ \\
modern $m^*(a)$ & 0.34809734, 0.85354081, 0.98074107, 0.99875760, 0.99999876 & 0.34809734, 0.85354081, 0.98074107, 0.99875760, 0.99999876 & $4.1\times10^{-15}$ \\
modern $p_b/\varepsilon$ at $a=100$ & 0.049773512 & 0.049773512 & $3.1\times10^{-15}$ \\
far future $m^*(x_0)$ & 0.34809734, 0.93957969, 0.99692358, 0.99999236, 1.00000000 & 0.34809734, 0.93957969, 0.99692358, 0.99999236, 1.00000000 & $4.1\times10^{-15}$ \\
Einstein run: $H_b(0)$; $a,b,H_a,H_b$ at $x_0=20$ & 0.0091437124, 1.5817387, 1.1975586, 0.011877999, 0.0074191078 & 0.0091437124, 1.5817387, 1.1975586, 0.011877999, 0.0074191078 & $2.6\times10^{-12}$ \\
Einstein–Lovelock run: $H_b(0)$; $a,b,H_a,H_b$ at $x_0=20$ & 0.014980704, 1.4773515, 1.2760841, 0.010175914, 0.0084887647 & 0.014980704, 1.4773515, 1.2760841, 0.010175914, 0.0084887647 & $2.1\times10^{-12}$ \\*
\end{longtable}}
The largest relative difference is $2.6\times10^{-12}$.

The checks of sections 15–16 of the notebook:

{\footnotesize
\begin{longtable}{>{\raggedright\arraybackslash}p{\dimexpr 0.0600\linewidth-2\tabcolsep\relax}>{\raggedright\arraybackslash}p{\dimexpr 0.7400\linewidth-2\tabcolsep\relax}>{\raggedright\arraybackslash}p{\dimexpr 0.0800\linewidth-2\tabcolsep\relax}>{\raggedright\arraybackslash}p{\dimexpr 0.1200\linewidth-2\tabcolsep\relax}}
\toprule
\textbf{\#} & \textbf{check} & \textbf{result} & \textbf{section} \\
\midrule\endfirsthead
\toprule
\textbf{\#} & \textbf{check} & \textbf{result} & \textbf{section} \\
\midrule\endhead
\endfoot
\bottomrule\endlastfoot
1 & \texttt{256 multi-indices in 9 ranks with 1, 8, 28, 56, 70, 56, 28, 8, 1 members} & True & 15.1 \\
2 & \texttt{every bilinear psibar Theta\textasciicircum{}K psi is real: Gamma17[0].Theta\textasciicircum{}K Hermitian for all 256 K} & True & 15.1 \\
3 & \texttt{Theta\_K.Theta\textasciicircum{}K == ID16 and Tr[Theta\_K.Theta\textasciicircum{}J] == 16 delta\_KJ} & True & 15.1 \\
4 & \texttt{Hodge duality: Theta\textasciicircum{}K.omega17 == (phase) Theta\textasciicircum{}(complement of K), |phase| = 1} & True & 15.1 \\
5 & \texttt{the couplings of section 8 are four of the nine: I\_0 == (psibar psi)\textasciicircum{}2, I\_1 == (psibar Gamma\textasciicircum{}A psi)(psibar Gamma\_A psi), I\_8 == -(psibar i chi17 psi)\textasciicircum{}2, I\_7 == -(axial invariant)} & True & 15.1 \\
6 & \texttt{local Lorentz invariance: I\_r(S psi, psibar S\textasciicircum{}-1) == I\_r(psi, psibar) for r = 0..8 (random finite Lorentz transformation), and S\textasciicircum{}-1 == Gamma17[0].S\textasciicircum{}dagger.Gamma17[0]} & True & 15.1 \\
7 & \texttt{Fierz matrix: closed form == (1/256) Sum\_\{K in r\} Tr[Theta\_K.Theta\_J.Theta\textasciicircum{}K.Theta\textasciicircum{}J] for the first and last J of every rank s} & True & 15.2 \\
8 & \texttt{F.F == IdentityMatrix[9]; eigenvalues: five +1 and four -1} & True & 15.2 \\
9 & \texttt{Fierz identity for commuting spinors (random): Sum\_\{K in r\} (a Theta\_K b)(c Theta\textasciicircum{}K d) == Sum\_s F\_rs Sum\_\{J in s\} (a Theta\_J d)(c Theta\textasciicircum{}J b), r = 0..8} & True & 15.2 \\
10 & \texttt{Hartree-Fock energy E\_H + E\_x == Sum\_s (G\_s/2) D\_s with G = (1 - Transpose[F]) g (random density matrix and couplings)} & True & 15.2 \\
11 & \texttt{uniform density matrix (rhoS ID16 + rho Gamma17[0])/16: D\_0 == rhoS\textasciicircum{}2, D\_1 == rho\textasciicircum{}2, D\_2 .. D\_8 == 0} & True & 15.2 \\
12 & \texttt{section 8 recovered: with g = (gS, gV, 0, 0, 0, 0, 0, -gA, -gP) the exchange energy of uniform matter is that of section 8.7} & True & 15.2 \\
13 & \texttt{the Hartree-Fock couplings of section 8.7: gS = -100000, gV = 10000 give G\_0 = -98750, G\_1 = 20000} & True & 15.2 \\
14 & \texttt{nonlinear field equation: dL\_self/d psibar\_a == -Sum\_r g\_r Sum\_K (psibar Theta\_K psi) (Theta\textasciicircum{}K psi)\_a} & True & 15.3 \\
15 & \texttt{chi17 commutes with every generator (Gamma17up[A].Gamma17up[B] - Gamma17up[B].Gamma17up[A])/4, hence with every Omega\_mu} & True & 15.4 \\
16 & \texttt{curved chiral equations: E\_L\textasciicircum{}dag.Gamma17[0].(i Gamma\textasciicircum{}mu D\_mu - m) psi == (i d0 + i(3Ha/2 + 2Hb) + tauO.d) phi\_L - m phi\_R, and the same for E\_R with -tauO.d} & True & 15.4 \\
17 & \texttt{flat limit a = b = 1, kF = kFsat: m* == 0.3480973401... and E/A - m == -0.0283214279... (section 8.7), to 10\textasciicircum{}-11} & True & 16.1 \\
18 & \texttt{flat limit: pa == pb == rho eps\_F - eps (Gibbs-Duhem), which is 0 at the saturation density} & True & 16.1 \\
19 & \texttt{energy-momentum conservation of the Kohn-Sham matter: eps' + 3Ha(eps + pa) + 4Hb(eps + pb) == 0 (a = 1 + 0.3x + 0.05x\textasciicircum{}2, b = 1/(1 + 0.2x), x0 = 0.7; relative 10\textasciicircum{}-8)} & True & 16.1 \\
20 & \texttt{primordial: rho x0 == kFsat\textasciicircum{}7/(105 Pi\textasciicircum{}4) at every x0, and m* decreases monotonically to below 0.05 at x0 = 10\textasciicircum{}-4} & True & 16.2 \\
21 & \texttt{modern (b = 1): eps a\textasciicircum{}3 tends to a constant and pa/eps -> 0 (pressureless in ordinary space)} & True & 16.2 \\
22 & \texttt{far future (de Sitter, H = 0.1): rho == rho0 e\textasciicircum{}\{-3 H x0\} and 1 - m* decays like rho} & True & 16.2 \\
23 & \texttt{the 11 and 44 Einstein-Lovelock equations are linear in (Ha', Hb'); the 00 equation contains neither (constraint)} & True & 16.3 \\
24 & \texttt{self-consistent evolution: the constraint Sum alpha\_k E\textasciicircum{}(k)0\_0 - Lambda == kappa eps is preserved to 10\textasciicircum{}-7 (relative) for Einstein and for Einstein-Lovelock (alpha2 = 20, alpha3 = 500)} & True & 16.3 \\*
\end{longtable}}
The Python program’s boolean checks:

{\footnotesize
\begin{longtable}{>{\raggedright\arraybackslash}p{\dimexpr 0.0600\linewidth-2\tabcolsep\relax}>{\raggedright\arraybackslash}p{\dimexpr 0.8200\linewidth-2\tabcolsep\relax}>{\raggedright\arraybackslash}p{\dimexpr 0.1200\linewidth-2\tabcolsep\relax}}
\toprule
\textbf{\#} & \textbf{check} & \textbf{result} \\
\midrule\endfirsthead
\toprule
\textbf{\#} & \textbf{check} & \textbf{result} \\
\midrule\endhead
\endfoot
\bottomrule\endlastfoot
1 & \texttt{256 products in 9 ranks: 1, 8, 28, 56, 70, 56, 28, 8, 1} & True \\
2 & \texttt{every bilinear psibar Theta\textasciicircum{}K psi is real: Gamma17[0] Theta\textasciicircum{}K Hermitian (Theta\textasciicircum{}K = i\textasciicircum{}((r(r-1)/2) mod 2) Gamma\textasciicircum{}K)} & True \\
3 & \texttt{Theta\_K Theta\textasciicircum{}K = 1\_16 and tr(Theta\_K Theta\textasciicircum{}J) = 16 delta\_KJ} & True \\
4 & \texttt{Fierz completeness: sum\_K (Theta\_K)\_cd (Theta\textasciicircum{}K)\_ab = 16 delta\_ad delta\_bc} & True \\
5 & \texttt{Hodge duality: Theta\textasciicircum{}K omega17 = +- Theta\textasciicircum{}(complement of K) (up to a phase)} & True \\
6 & \texttt{stage-2 channels: I\_0 = (psibar psi)\textasciicircum{}2, I\_1 = (psibar Gamma\textasciicircum{}A psi)(psibar Gamma\_A psi), I\_8 = -(psibar i chi psi)\textasciicircum{}2, I\_7 = -(axial)} & True \\
7 & \texttt{local Lorentz invariance: I\_r(S psi, psibar S\textasciicircum{}-1) = I\_r(psi, psibar), r = 0..8 (random finite boost-rotation)} & True \\
8 & \texttt{Fierz matrix F\_rs = (1/256) sum\_\{K in r\} tr(Theta\_K Theta\_J Theta\textasciicircum{}K Theta\textasciicircum{}J), independent of the representative J of rank s} & True \\
9 & \texttt{Fierz matrix closed form F\_rs = (1/16) sum\_l (-1)\textasciicircum{}(rs-l) C(s,l) C(8-s,r-l)} & True \\
10 & \texttt{F\textasciicircum{}2 = 1 (the Fierz transformation is an involution)} & True \\
11 & \texttt{Fierz identity for commuting spinors: sum\_\{K in r\} (a Theta\_K b)(c Theta\textasciicircum{}K d) = sum\_s F\_rs sum\_\{J in s\} (a Theta\_J d)(c Theta\textasciicircum{}J b)} & True \\
12 & \texttt{Hartree-Fock energy E\_H + E\_x = sum\_s (G\_s/2) D\_s with G = (1 - F\textasciicircum{}T) g (random density matrix, random couplings)} & True \\
13 & \texttt{uniform matter R = (rho\_S + rho Gamma17[0])/16: only D\_0 = rho\_S\textasciicircum{}2, D\_1 = rho\textasciicircum{}2; with g\_0 = gS, g\_1 = gV, g\_7 = -gA, g\_8 = -gP the exchange energy is that of stage 2 (section 8.7)} & True \\
14 & \texttt{uniform density matrix: D\_0 = rho\_S\textasciicircum{}2, D\_1 = rho\textasciicircum{}2, D\_2..D\_8 = 0} & True \\
15 & \texttt{stage-2 couplings gS = -100000, gV = 10000: G\_0 = -98750, G\_1 = 20000 (the HF couplings of section 8.7)} & True \\
16 & \texttt{nonlinear field equation: dL\_int/d psibar\_a = -sum\_r g\_r sum\_K (psibar Theta\_K psi)(Theta\textasciicircum{}K psi)\_a} & True \\
17 & \texttt{the chirality chi17 commutes with every generator Sigma\textasciicircum{}AB, hence with every spin connection Omega\_mu} & True \\
18 & \texttt{curved chiral equations: E\_L\textasciicircum{}dag Gamma0 (i Gamma\textasciicircum{}mu D\_mu - m) psi = (i d0 + i(3Ha/2 + 2Hb) + tauO.d/(a,b)) phi\_L - m phi\_R} & True \\
19 & \texttt{curved chiral equations: E\_R\textasciicircum{}dag Gamma0 (i Gamma\textasciicircum{}mu D\_mu - m) psi = (i d0 + i(3Ha/2 + 2Hb) - tauO.d/(a,b)) phi\_R - m phi\_L} & True \\
20 & \texttt{flat limit a = b = 1, kF = kF0 (stage-2 saturation): m* = 0.3480973401..., E/A - m = -0.0283214279...} & True \\
21 & \texttt{flat limit: isotropic pressure pa = pb = rho eps\_F - eps (Gibbs-Duhem, eps\_F = V0 + sqrt(kF\textasciicircum{}2 + m*\textasciicircum{}2))} & True \\
22 & \texttt{energy-momentum conservation eps' + 3Ha(eps + pa) + 4Hb(eps + pb) = 0 for the KS matter (arbitrary a, b; relative residual)} & True \\
23 & \texttt{primordial: rho(x0) x0 = kF0\textasciicircum{}7/(105 pi\textasciicircum{}4) (density grows like 1/x0)} & True \\
24 & \texttt{primordial: m* decreases toward 0 as x0 -> 0 (blue-shifted ordinary momenta, attractive scalar coupling)} & True \\
25 & \texttt{modern: eps a\textasciicircum{}3 -> constant and pa/eps -> 0 (dirac16complex behaves as pressureless matter in ordinary space)} & True \\
26 & \texttt{far future: rho = rho0 e\textasciicircum{}\{-3 H x0\} and 1 - m* -> 0 like rho} & True \\
27 & \texttt{the 00 closed forms of E\textasciicircum{}(1), E\textasciicircum{}(2), E\textasciicircum{}(3) (pair-set algorithm, tools/lovelock71.py) agree with the closed forms of section 12.6 of 7+1.nb} & True \\
28 & \texttt{Einstein-Lovelock equations are linear in (Ha', Hb') (quasi-linear second-order system)} & True \\
29 & \texttt{the constraint Sum alpha\_k E\textasciicircum{}(k)0\_0 - Lambda = kappa eps is preserved by the evolution (relative violation < 1e-7) in both cases} & True \\*
\end{longtable}}
\Needspace{12\baselineskip}\section*{10. How to reproduce}\addcontentsline{toc}{section}{10. How to reproduce}
\begin{enumerate}
\item \texttt{python3 tools/make\_7p1\_nb.py} writes 7+1.nb (unevaluated).
\item Evaluate 7+1.nb in a fresh kernel of Wolfram 15.0.1 (Evaluation \texorpdfstring{\ensuremath{\triangleright}}{▸} Quit Kernel, then Evaluation \texorpdfstring{\ensuremath{\triangleright}}{▸} Evaluate Notebook; or \texttt{Get["D:/Developer/github/7+1/tools/run\_eval.wl"]}) and save it; the last section summarizes all checks.
\item \texttt{python3 tools/ref\_ks\_gravity.py} repeats the checks of this document independently (exit code 0 if all pass).
\item \texttt{python3 tools/make\_docs.py} regenerates this document, its symbol dictionary and its provenance record.
\end{enumerate}
\end{document}
